This section describes the realised artifact: the components that implement the architecture of
Section~\ref{sec:proposed_system}, and the technologies each is built on. It reports what exists,
not how well it performs. \textbf{No integration testing across the complete pipeline and no
evaluation have been carried out}; the protocol for the latter is given in
Section~\ref{sec:protocol}, and the consequences for what may be claimed are stated in
Section~\ref{sec:limitations}.

\subsection{Component Structure}

The system is implemented as the four deployable components listed in
Table~\ref{tab:components}, communicating through a Firebase Realtime Database. The database
decouples them: the edge node does not need to know whether a dashboard is connected, and the
vision client and backend can be restarted independently of the hardware.

\begin{table}[htbp]
\centering
\caption{Implemented components and their technologies}
\label{tab:components}
\begin{tabular}{@{}llp{6.2cm}@{}}
\toprule
\textbf{Component} & \textbf{Technology} & \textbf{Responsibility} \\
\midrule
Edge firmware
& C++ / Arduino, PlatformIO
& Camera streaming server, ultrasonic and photoresistor sampling, two-colour indicator driven from
the database, publication of raw telemetry \\
\addlinespace
Vision client
& Python, MediaPipe, OpenCV
& Face mesh extraction, head pose by \texttt{solvePnP}, Eye Aspect Ratio and sliding-window blink
rate, interpupillary distance estimate, per-flag persistence timers \\
\addlinespace
Backend service
& Python, FastAPI, SQLite
& Session state machine, flag detection and qualification, severity assignment, notification
governor, LLM invocation, session history \\
\addlinespace
Dashboard
& HTML, CSS, vanilla JavaScript
& Live telemetry, guided calibration, notification surface with governor actions, calibration and
threshold configuration \\
\bottomrule
\end{tabular}
\end{table}

\subsection{Design Choices Worth Recording}

Three implementation decisions follow from the architecture rather than from convenience, and are
recorded here because they constrain how the system can be deployed and extended.

\textbf{A single configuration source.} Every threshold, hold time, severity rule and governor
parameter in Table~\ref{tab:parameters} is read from one configuration file on each request, and
the calibration values are editable from the dashboard at runtime. This exists so that the values
described in this paper are demonstrably the values the system uses, and so that a deployment can
be retuned for a different user or workspace without redeployment.

\textbf{The governor is the only writer of user-facing state.} Session persistence, indicator
state, history, and LLM invocation are all driven from the session engine rather than from
independent timers. Without this, a page reload or a second open dashboard would reset cooldowns
and re-raise an alert the user had already dismissed --- the failure mode that makes notification
management ineffective in practice.

\textbf{Deterministic fallback for the recommendation layer.} When the language model is
unreachable, a per-flag phrase table supplies the recommendation and the record is labelled with
the fallback model name. This keeps the system usable offline, and it also means generated and
templated output can later be compared directly on identical detected states, which is the basis
of the comparison described in Section~\ref{sec:protocol}.

\subsection{Illuminance Conversion}

The photoresistor is sampled as a raw 12-bit value and converted to lux in the backend rather than
on the device, so that the conversion has a single definition. Two properties of that conversion
are worth stating because they bound what the illuminance channel can report.

First, the divider is inverted to the photoresistor's own resistance in normalised form, in which
the supply voltage cancels; supply voltage is therefore not a parameter, which removes a source of
error given that the regulated rail is not exactly nominal.

Second, the conversion is bounded at both ends. Above the point at which the analogue-to-digital
converter saturates, the photoresistor still has finite resistance, so a saturated sample means
``at least this bright'' rather than a specific value. The implementation returns a range status
alongside the value and reports such samples as a lower bound. Treating saturation as zero
illuminance --- which a naive divider inversion does, since it computes zero resistance at full
scale --- would invert the reading and raise a low-light warning under the brightest conditions the
sensor can register.

\subsection{Status}

All components described in Section~\ref{sec:proposed_system} are implemented, and the individual
processing stages have been exercised during development. What has \emph{not} been done is
end-to-end integration testing of the complete pipeline under realistic conditions, and no
measurement of accuracy, latency, interruption load, or user outcome is reported anywhere in this
paper. Establishing those is the immediate next step and is specified in Section~\ref{sec:protocol}.
